The coordination of the "Virtual-Europe-Path" project is led by Eötvös Loránd University (ELTE), representing the Central and Eastern European (CEE) strategic axis. As the coordinating institution, ELTE provides the administrative, legal, and academic framework necessary to manage a large-scale, multi-national EU consortium.

\textbf{Dr. Krémer Aliz - Project Director and Managing Director}
Dr. Krémer Aliz serves as the primary coordinator and legal representative of the consortium. In this role, she oversees the alignment of all nine Work Packages (WPs) and ensures the timely delivery of project milestones.

\textbf{Professional Profile and Expertise}
\begin{itemize}
    \item \textbf{Strategic Management}: Extensive experience in leading interdisciplinary research teams at the intersection of Digital Humanities and Cultural Heritage.
    \item \textbf{EU Project Coordination}: Specialized in the administrative and financial management of Horizon Europe and Erasmus+ frameworks, ensuring compliance with Open Science and GDPR principles.
    \item \textbf{Regional Expertise}: A leading advocate for the digital preservation of endangered CEE heritage, specifically focusing on the architectural vulnerability of Transylvanian wooden churches and manor houses.
\end{itemize}

\textbf{Key Responsibilities within the Project}
\begin{itemize}
    \item \textbf{Consortium Leadership}: Acting as the main point of contact between the European Commission and the partner institutions (Arts et Métiers, Sorbonne, Sapienza, NKUA, ELTE, and Springer Nature).
    \item \textbf{Quality and Risk Oversight}: Directing the development of the Quality Assurance (QA) and Risk Management Plan to mitigate technical or logistical delays in field scanning.
    \item \textbf{Financial Governance}: Supervising the budget allocation and financial reporting across all partner regions (Mediterranean, Atlantic, and CEE).
    \item \textbf{Dissemination Strategy}: Working closely with Springer Nature to ensure that the project's high-fidelity data and scientific outputs achieve maximum global impact.
\end{itemize}

\textbf{Contact Information}
